\documentclass[aps,onecolumn,nofootinbib,pra]{article}

\usepackage{../../spec_files/arxiv}
\input{../../spec_files/standard_preamble.tex}

\usepackage{import}
\subimport{../../macros/}{all-macros}

\usepackage{minted}

\newcommand{\satsemiring}{(\{0,1\},\lor,\land,0,1)}
\newcommand{\gensemiring}{(S,\oplus,\otimes,0,1)}
\newcommand{\srcontractionwrt}[2]{\langle#1\rangle^{#2}}
\newcommand{\srcontractionofwrt}[3]{\langle#1\rangle^{#3}_{#2}}

\begin{document}
    \title{Semiring Tensor Networks}

    \maketitle
    \date{\today}


    \section{Semiring}

    \begin{definition}
        A semiring is a tuple $\gensemiring$, where $S$ is a set, $0,1\in S$, and
        \begin{align*}
            \oplus \defcols S \times S \rightarrow S \quad, \quad \otimes \defcols S \times S \rightarrow S \, ,
        \end{align*}
        such that the following properties hold:
        \begin{itemize}
            \item[(i)] $(S,\oplus,0)$ is a commutative monoid.
            \item[(ii)] $(S,\otimes,1)$ is a commutative monoid.
            \item[(iii)] $\otimes$ distributes over $\oplus$, that is for any $s_0,s_1,s_2\in S$
            \begin{align*}
                s_0 \otimes (s_1 \oplus s_2)
                = (s_0 \otimes s_1) \oplus (s_0 \otimes s_2) \, .
            \end{align*}
            \item[(iv)] $0$ annihilates $\otimes$, that is (check!)
            \begin{align*}
                0 \otimes (s_0\otimes s_1) = 0 \, .
            \end{align*}
        \end{itemize}
    \end{definition}

    We understand $S$ as the set of coordinates in a tensor.

    \begin{definition}
        A tensor with coordinates in $S$ of order $\catorder\in\nn$ and leg dimensions $\catdimof{[\catorder]}=(\catdimof{0},\ldots,\catdimof{\catorder-1})$, where $\catdimof{\catenumerator}\in\nn$ for $\catenumeratorin$, is a map
        \begin{align*}
            \hypercore \defcols \facstates \rightarrow S \, .
        \end{align*}
        A tensor network with coordinates in $S$ is a hypergraph $\graph=(\nodes,\edges)$, where each node is decorated by a dimension $\catdimof{\node}$ and each edge by a tensor of .
    \end{definition}


    \section{Contractions}

    \begin{definition}
        Given a semiring $\gensemiring$ and a tensor network $\extnetasset$ with coordinates in $S$, the contraction of the tensor network keeping the variables to a subset $\arbset\subset\nodes$ open is the tensor
        \begin{align*}
            \srcontractionofwrt{
                \extnetasset
            }{\catvariableof{\arbset}}{\gensemiring}
        \end{align*}
        with coordinates at indices $\catindexof{\arbset}$ by
        \begin{align*}
            \srcontractionofwrt{
                \extnetasset
            }{\catvariableof{\arbset}}{\gensemiring}
            = \bigoplus_{\catindexof{\nodes\backslash\arbset}} \bigotimes_{\edgein} \hypercoreofat{\edge}{\indexedcatvariableof{\edge}} \, .
        \end{align*}
    \end{definition}

    The distributive property $(iii)$ of the semiring ensures, that we can compute contractions by sub-contractions.

    \begin{theorem}
        Given a disjoint partition of $\edges$ into $\edgesof{0},\edgesof{1}$, and a subset $\arbset\subset\nodes$, we define the sets $\nodesof{0}$, $\nodesof{1}$ and $\nodesof{0\cup1}$ as the nodes appearing in $\edgesof{0}$, respectively $\edgesof{1}$ and in $\edgesof{0}$ and $\edgesof{1}$.
        We then have
        \begin{align*}
            \srcontractionofwrt{
                \extnetasset
            }{\catvariableof{\arbset}}{\gensemiring}
            =
            \srcontractionofwrt{
                \srcontractionofwrt{
                    \{\hypercoreofat{\edge}{\catvariableof{\edge}} \wcols \edge\in\edgesof{0}\}
                }{\catvariableof{\nodesof{0}\cap\arbset},\catvariableof{\nodesof{0\cup1}}}{\gensemiring},
                \srcontractionofwrt{
                    \{\hypercoreofat{\edge}{\catvariableof{\edge}} \wcols \edge\in\edgesof{1}\}
                }{\catvariableof{\nodesof{1}\cap\arbset},\catvariableof{\nodesof{0\cup1}}}{\gensemiring}
            }{\catvariableof{\arbset}}{\gensemiring} \, .
        \end{align*}
    \end{theorem}
    \begin{proof}
        We use the commutative properties to split the summations and the distributive property to move the $\bigoplus$ to closed variables only in $\nodesof{0}$ respectively $\nodesof{1}$ to the first respectively second sub-contraction.
    \end{proof}

    \section{Examples}

    \subsection{Standard TN}

    Standard TN use the semiring
    \begin{align*}
        (\rr,+,\cdot,0,1) \, .
    \end{align*}

    \subsection{Constraint Satisfaction}

    We take $S=\{0,1\}$ and interpret $0$ as $\falsesymbol$, $1$ as $\truesymbol$.
    We define the disjunction and conjunction operations by
    \begin{align*}
        \lor \left[\catvariableof{0},\catvariableof{1}\right]
        = \coloredmatrixof{0 & 1 \\ 1 & 1}{\catvariableof{0},\catvariableof{1}}
    \end{align*}
    and
    \begin{align*}
        \land \left[\catvariableof{0},\catvariableof{1}\right]
        = \coloredmatrixof{0 & 0 \\ 0 & 1}{\catvariableof{0},\catvariableof{1}} \, .
    \end{align*}

    Tensors over $\{0,1\}$ are interpreted as constraints, and their tensor networks as constraint networks.

    The satisfiability of a constraint network $\extnetasset$ is the contraction
    \begin{align*}
        \bigvee_{\catindexof{\nodes}} \bigwedge_{\edgein} \hypercoreofat{\edge}{\indexedcatvariableof{\edge}}
        = \srcontractionofwrt{
            \extnetasset
        }{\varnothing}{\satsemiring} \, .
    \end{align*}

%    This is the contraction
%    \begin{align*}
%        \langle \extnetasset\rangle^{(\{0,1\},\lor,\land,0,1)}
%    \end{align*}

    \subsection{Maximum calculus}

    The maximal coordinate of a standard TN is the contraction wrt the semiring
    \begin{align*}
        (\rr,\mathrm{max},\cdot,0,1) \, .
    \end{align*}

    Analogously the minimal coordinate by
    \begin{align*}
        (\rr,\mathrm{min},\cdot,0,1) \, .
    \end{align*}

    These schemes only work, when taking the minimal coordiante keeping all variables open.
    When searching for the maximal/minimal coordiante of a contraction with some variables closed, we need to first coarse grain the tensor network.

    These contractions are used in encoding-decoding schemes \cite{mezard_information_2009,mackay_information_2003}.



    \bibliographystyle{plainnat}
    \bibliography{../../references.bib}


\end{document}